\documentclass[11pt]{article}
\usepackage[margin=1in]{geometry}
\usepackage{hyperref}
\usepackage{listings}
\usepackage{xcolor}
\usepackage{booktabs}
\usepackage{enumitem}
\usepackage{amsmath}

\emergencystretch=3em
\tolerance=2000
\sloppy

\lstset{
  basicstyle=\ttfamily\small,
  breaklines=true,
  columns=fullflexible,
  frame=single,
  backgroundcolor=\color{gray!10},
  showstringspaces=false
}

\title{User-Facing Descriptions of the URI OPeNDAP Datasets\\
       \large grep-dap test case, Prompt \#5}
\author{Compiled for P. Cornillon}
\date{2026-05-16}

\begin{document}
\maketitle

\begin{abstract}
Short, accessible descriptions of the two publicly readable
datasets on \url{https://sst-aqua.gso.uri.edu/opendap/}, written
for a working scientist deciding whether the data fits a need and
how to pull a sub-set. Detailed agent-facing instructions are
shipped as the two companion files
\texttt{docs/usage\_sst.md} and
\texttt{docs/usage\_gradient\_SST.md}; this document summarizes
what those files cover.
\end{abstract}

\tableofcontents

\section{Dataset 1: \texttt{SST\_Orbits/}}\label{sec:ds1}

\subsection{At a glance}

\begin{description}[leftmargin=2em,style=nextline,topsep=2pt]
  \item[What it is.]
    Per-orbit, L2-swath sea-surface temperature from MODIS aboard
    NASA's Aqua satellite, reprocessed by URI. Each file is one
    complete Aqua orbit ($\sim$100 minutes, $\sim$40\,000
    along-track lines of 1\,354 across-track 1-km pixels).
  \item[Coverage.]
    2002 through 2024; $\sim$14 orbits/day; $\sim$1.2$\times 10^5$
    files in total, organized by year and month.
  \item[Native resolution.] 1 km at nadir, growing to $\sim$5 km at
    the swath edges. MODIS swath width 2\,330 km.
  \item[Format.] NetCDF-4, served via OPeNDAP Hyrax (DAP4 preferred).
  \item[What you get per file.] SST, an SST gradient field (eastward
    and northward components, in $^\circ$C/km), a quality-flag field,
    a refined-mask field, pixel-level latitude/longitude, per-scan
    nadir and swath-edge geolocation, the orbit start time, and a
    side-group of co-located AMSR-E microwave fields (only meaningful
    while AMSR-E was operational, June 2002 -- October 2011).
  \item[Producer.]
    P.\ Cornillon, URI Graduate School of Oceanography, ``SST Fronts''
    project (\url{http://www.sstfronts.org}).
\end{description}

\subsection{Who it is for}

This is the right dataset if you want \emph{per-orbit, swath-level}
MODIS-Aqua SST and gradient values --- for example, to detect SST
fronts in a single overpass, to validate against in-situ
measurements, to study individual mesoscale features, or to build
your own gridded composites. The orbit structure is preserved, so
the scan geometry (along-scan / along-track, nadir / edge) is
recoverable, and the file's quality flags and refined mask let a
careful user filter pixels by data quality.

It is \emph{not} the right dataset if you want a Level-3 or
Level-4 SST analysis (gap-filled, gridded to a regular lat/lon
mesh). For that, use a product such as JPL MUR
or NOAA/NCEI's daily L4 (not provided here in publicly readable
form).

\subsection{How to get a sub-set}

The OPeNDAP server allows array-index sub-setting via standard
DAP4 hyperslab requests. The cheapest way is to first pull
\texttt{nadir\_latitude} and \texttt{nadir\_longitude} (both of
length \texttt{ny}, so $\sim$40\,000 floats each --- a few hundred
kilobytes), bracket the line indices you want, then request the
science variables over that index range:

\begin{lstlisting}[language=Python]
from pydap.client import open_url
import numpy as np

URL = ("https://sst-aqua.gso.uri.edu/opendap/SST_Orbits/2024/01/"
       "AQUA_MODIS_orbit_115220_20240101T011558_L2_SST-URI_24-2.nc4")
ds = open_url(URL, protocol="dap4")

# Find lines near a specific nadir latitude
nlat = np.asarray(ds["nadir_latitude"][:].data, dtype="f8")
i0, i1 = np.searchsorted(nlat, [-5, +5])  # tropical band of this orbit

# Pull a small SST patch in the middle of the swath
sst   = ds["SST_In"][i0:i1, 600:700]
scale = float(ds["SST_In"].attributes["scale_factor"])
fill  = float(ds["SST_In"].attributes["_FillValue"])
data  = np.asarray(sst.data, dtype="f8")
data[data == fill] = np.nan
data *= scale  # now in degree_Celsius
\end{lstlisting}

\noindent
Full agent-facing instructions, including all 44 variables in the
three groups, code recipes for bounding-box sub-setting and
multi-orbit aggregation, and the metadata pitfalls that catch
new users (e.g.\ \texttt{units="C"} not being UDUNITS, AMSR-E
fields being $1\times 1$ placeholders after October 2011, etc.):
see \texttt{docs/usage\_sst.md}.

\subsection{Honest limitations}

\begin{itemize}
  \item Some variables declare units (\texttt{"C"}, \texttt{"C/km"},
    \texttt{seconds}) that are not strictly UDUNITS-compatible; a
    strict CF reader may reject them. They are unambiguous in
    context.
  \item The quality-flag attribute uses \texttt{flag\_masks=0},
    which is malformed. Interpret \texttt{qual\_sst} as integer
    values $0$\,=BEST, $1$\,=GOOD, $2$\,=QUESTIONABLE, $3$\,=BAD,
    $4$\,=NOTPROCESSED.
  \item The co-located AMSR-E variables are $1\times 1$
    placeholders in files dated after October 2011.
\end{itemize}

\section{Dataset 2: \texttt{gradients\_by\_period/}}\label{sec:ds2}

\subsection{At a glance}

\begin{description}[leftmargin=2em,style=nextline,topsep=2pt]
  \item[What it is.]
    Monthly statistics on a $1^\circ \times 1^\circ$ global grid of
    SST and SST-gradient values originally computed per-pixel from
    the URI L2 MODIS-Aqua orbits (Dataset~1). Each cell carries
    sums, sums of squares, and pixel counts --- enough to recover
    monthly means, variances, and 1$^\circ$-cell distributions.
  \item[Coverage.]
    January 2003 through December 2024 (270 monthly files);
    global; 1$^\circ$ bins; day-side and night-side separately.
  \item[Spatial scale of the underlying gradient.]
    $\sim$5 km (operator support of the differentiation step
    applied to the L2 SST); not the 1$^\circ$ bin size.
  \item[Format.] NetCDF, served via OPeNDAP Hyrax. \emph{The
    files declare no semantic metadata at all} --- no units, no
    fill value, no coordinate variables, no long names, no globals.
  \item[What you get per file.] 34 \texttt{Float64} variables of
    shape \texttt{(lon=360, lat=180)}: day/night pixel counts (three
    flavours), sums of SST and SST squared, sums of two coordinate
    forms of the gradient (geographic eastward / northward, and
    swath-relative along-scan / along-track) and their magnitudes,
    with paired \texttt{\_squared} variables. No time variable; the
    file's month is encoded in the filename only.
  \item[Producer.]
    Same as Dataset~1 (inferred --- the gradient algorithm and unit
    semantics descend directly from the orbit branch).
\end{description}

\subsection{Who it is for}

This is the right dataset if you want \emph{monthly, gridded}
summaries of SST and SST-gradient statistics across many years,
for example to compute climatologies of SST fronts, to look at
mesoscale activity patterns by season, to compare hemispheres or
basins, or to construct training labels for ML models that work
on gridded fields. The bookkeeping ($\Sigma$, $\Sigma^2$, $N$ per
cell) means you can compute not just per-cell means but also
variances, and you can re-aggregate across months by simply
summing.

It is \emph{not} the right dataset if you need the original
per-pixel gradient values; for those, use the source orbit files
in \texttt{SST\_Orbits/} (Dataset~1). The 1$^\circ$ binning hides
the actual $\sim$5 km gradient measurement scale, so do not
advertise the data as ``$1^\circ$ SST gradient.''

\subsection{How to get a sub-set}

The grid is small ($360\times 180$ Float64 = 540 kB per variable),
so often the cheapest thing is to pull the whole field for one
month. For a bounding-box sub-set:

\begin{lstlisting}[language=Python]
from pydap.client import open_url
import numpy as np

URL = ("https://sst-aqua.gso.uri.edu/opendap/gradients_by_period/"
       "augmented_monthly_stats_07_2020.nc")
ds = open_url(URL, protocol="dap4")

# Inferred grid: cell-centered, south-to-north, -180..+180
lonidx = lambda lon: int(round(lon + 179.5))
latidx = lambda lat: int(round(lat +  89.5))

i_lon0, i_lon1 = lonidx(-100), lonidx(-50) + 1   # 100W..50W (Atlantic)
i_lat0, i_lat1 = latidx( 20),  latidx( 50) + 1   # 20N..50N

N = np.asarray(ds["day_pixel_count"][i_lon0:i_lon1, i_lat0:i_lat1].data, dtype="f8")
S = np.asarray(ds["day_sum_SST"   ][i_lon0:i_lon1, i_lat0:i_lat1].data, dtype="f8")
mean_sst = np.where(N > 0, S / N, np.nan)        # degree Celsius
\end{lstlisting}

\noindent
Full agent-facing instructions --- including the verified grid
orientation, the per-variable inferred units, the divisor rule
(which pixel count to pair with which sum), a worked recipe for
recovering per-cell variance, and the list of every inference that
the file does not declare --- are in
\texttt{docs/usage\_gradient\_SST.md}.

\subsection{Honest limitations}

\begin{itemize}
  \item The file declares no semantic metadata. Everything described
    above is inferred from the variable names, from the companion
    orbit branch, and from direct data probes. The inferences are
    high-confidence for the structural items (grid, signs, fill,
    units, divisors); medium-confidence for the day/night
    interpretation (Aqua ascending vs. descending); and not at all
    inferable for the file's exact processing history. See
    \texttt{docs/uri\_test\_case\_3.tex} for the full audit.
  \item NaN appears in some sum variables even where the pixel
    count is positive --- an undeclared fill value. Always mask
    before dividing.
  \item Multiple flavours of gradient and pixel count: the
    swath-relative \texttt{grad\_as}/\texttt{grad\_at} should be
    divided by \texttt{as\_pixel\_count} and \texttt{at\_pixel\_count}
    respectively, not by \texttt{pixel\_count}; conversely the
    geographic \texttt{eastward}/\texttt{northward} gradients
    \emph{are} paired with \texttt{pixel\_count}. The mismatch
    catches most new users.
\end{itemize}

\section{Where the agent-facing files live}

\begin{itemize}
  \item \texttt{docs/usage\_sst.md} --- complete usage guide for
    \texttt{SST\_Orbits/}, including every variable, every pitfall,
    and self-contained Python recipes.
  \item \texttt{docs/usage\_gradient\_SST.md} --- complete usage
    guide for \texttt{gradients\_by\_period/}, including the
    inferred coordinate grid, units, divisors, and the full
    list of items the audit could not infer.
\end{itemize}

Both \texttt{.md} files are self-contained: an AI agent can be
handed one (or both) and use it as the sole reference for accessing
the data, without needing the broader \texttt{grep-dap} report
context.

\end{document}
